\subsection{Common native clock and wave observations}
On the same prepared energy-lapse geometry, every native transition has
frequency $\omega_k(x)=e^{q_x}(E_k-E_0)$. Complete eight-output Ramsey
instruments at two sites give 64 joint records and the common ratio
$\omega_k(a)/\omega_k(b)=e^{q_a-q_b}$ for all seven excited levels.
The source reference and probe preparation work are retained. Geometry is
fixed during the probes, without an autonomous backreaction assertion.

Declare the scalar coefficient density
$C=\operatorname{diag}(-e^{-4q},1,1,1)$. Its metric dictionary gives
$g=\operatorname{diag}(-e^{2q},e^{-2q},e^{-2q},e^{-2q})$, with that same
lapse. All vacuum and one-quantum modes of the finite field have
$K_{\rm wave}=e^{2q}D^TD e^{2q}$; complete site-frame observations retain
measurement work and independent evolution controls. The metric ray time
$\sum_j e^{-2q_j}$ has a certified positive delay, distinct from a claimed
lattice group-delay limit.

Independent first-order curvature gives $G_{00}^{(1)}=2\Delta q$ and
$G_{ij}^{(1)}=0$. Native energy stationarity therefore bounds
$G_{00}^{(1)}-2\rho_q/\kappa$ by the actual geometry residual, with
$\rho_q=e^qE$ defined from matter energy. This uses no installed Einstein
source equation. The isotropic spatial gluing is a declared choice:
a changed spatial metric can preserve the clock and force laws while
violating the exterior curvature comparison. Physical units, full
covariant stress and nonlinear Einstein dynamics are separate conditions.
